% Figura 3: camadas da BFS no tabuleiro de exemplo e um caminho mínimo de C a S.
%
% Cada casa é dada por um par "dígito/distância": o dígito é o do tabuleiro (x marca as
% casas C e S, desenhadas à parte) e a distância é a do mapa impresso por
% MedicoesHiperpulos (ver Trabalho2/medicoes-referencia.txt). A cor da casa vem da distância.
% Para mudar as cores, edite os \colorlet; para mover a legenda, mude o valor de \xl.
\begin{tikzpicture}[x=0.82cm,y=0.82cm,font=\footnotesize]
  \colorlet{nivel0}{black}
  \colorlet{nivel1}{blue!10}
  \colorlet{nivel2}{blue!25}
  \colorlet{nivel3}{blue!40}
  \colorlet{nivel4}{blue!58}
  \colorlet{trilha}{orange!85!black}

  % tabuleiro: uma lista por linha, de cima para baixo
  \foreach \lin [count=\r from 0] in {
    {4/1,2/2,8/3,6/2,3/3,0/4,5/1,9/2,9/3,3/4},
    {6/2,7/3,1/2,2/3,6/4,5/3,1/2,4/3,x/0,1/3},
    {3/1,5/2,4/3,1/2,5/3,5/2,9/1,2/2,3/3,8/4},
    {6/2,3/3,4/2,4/3,3/2,6/3,1/2,2/1,3/4,6/1},
    {0/3,4/4,1/3,9/2,8/3,6/2,7/3,8/4,0/3,2/2},
    {4/2,2/3,7/2,1/3,6/2,1/3,3/2,3/3,5/4,0/3},
    {x/3,6/2,6/3,8/2,9/3,6/2,2/3,6/2,8/3,5/2},
    {6/2,2/3,4/2,4/3,0/2,5/3,1/2,1/3,3/2,9/3},
    {3/3,7/2,1/3,4/4,7/3,0/2,2/3,0/2,3/3,0/4},
    {0/2,0/3,8/4,4/3,8/2,2/3,5/4,5/1,8/4,7/1}%
  } {
    \foreach \dig/\dis [count=\c from 0] in \lin {
      \draw[fill=nivel\dis] (\c,{-\r-1}) rectangle ++(1,1);
      \ifx\dig x\else \node[black!75] at ({\c+0.5},{-\r-0.5}) {\dig}; \fi
    }
  }

  % casa C (partida) e casa S (saída)
  \node[white,font=\footnotesize\bfseries] at (8.5,-1.5) {C};
  \draw[trilha,very thick] (0,-7) rectangle ++(1,1);
  \node[font=\footnotesize\bfseries] at (0.5,-6.5) {S};

  % numeração de linhas e colunas
  \foreach \i in {0,...,9} {
    \node[gray!50!black] at ({\i+0.5},0.4) {\i};
    \node[gray!50!black] at (-0.4,{-\i-0.5}) {\i};
  }

  % caminho mínimo: C -> (2,6) -> (2,7) -> S
  \tikzset{seta/.style={trilha,very thick,-{Stealth[length=2.2mm]},shorten >=3pt,shorten <=3pt},
           num/.style={circle,fill=white,draw=trilha,inner sep=1pt,font=\scriptsize\bfseries,text=trilha}}
  \draw[trilha,very thick] (6,-3) rectangle ++(2,1);  % casas (2,6) e (2,7), para o digito ficar legivel
  \draw[seta] (8.5,-1.5) -- (6.5,-2.5) node[num,midway] {1};
  \draw[seta,shorten >=6pt] (6.5,-2.5) -- (7.5,-2.5) node[num,midway,above=3pt] {2};
  % o pulo 3 sai pela borda direita (coluna 10) e entra pela esquerda (coluna 0)
  \draw[trilha,very thick,shorten <=3pt] (7.5,-2.5) -- (10,-5.833) node[num,pos=0.45] {3};
  \draw[seta,shorten <=0pt,shorten >=9pt] (0,-5.833) -- (0.5,-6.5);
  \fill[trilha] (10,-5.833) circle (0.1);   % ponto onde o pulo sai do tabuleiro
  \fill[trilha] (0,-5.833) circle (0.1);    % ponto onde ele volta a entrar

  % legenda
  \def\xl{10.7}
  \node[anchor=west,align=left] at (\xl,-0.6) {\textbf{Distância até C}\\(em pulos)};
  \foreach \k in {0,...,4} {
    \draw[fill=nivel\k] (\xl,{-2.2-0.85*\k}) rectangle ++(0.7,0.7);
    \node[anchor=west] at ({\xl+0.9},{-1.85-0.85*\k}) {\k};
  }
  \node[anchor=west,align=left,black!75] at (\xl,-7.0) {Número na casa:\\dígito do tabuleiro};
  \node[num] at ({\xl+0.25},-8.3) {1};
  \node[anchor=west,align=left,black!75] at ({\xl+0.6},-8.3) {ordem dos pulos};
\end{tikzpicture}
